\documentclass[prb,12pt, nofootinbib]{revtex4-2}
%fonts
% Palatino for main text and math
\usepackage[osf,sc]{mathpazo}

% Helvetica for sans serif
% (scaled to match size of Palatino)
\usepackage[scaled=0.90]{helvet}

% Bera Mono for monospaced
% (scaled to match size of Palatino)
\usepackage[scaled=0.85]{beramono}
\usepackage{amsmath, amssymb,physics,amsfonts,amsthm}
\usepackage{enumitem}
\usepackage{cancel}
\usepackage[most]{tcolorbox}
\usepackage{booktabs}
\usepackage{subcaption}
\captionsetup[subfigure]{% changed <<<<<<<<<<<<<<<<
	singlelinecheck = false,
	justification=raggedright, 
	margin = {-3ex, 0mm}, % make margin font size dependent
}
\usepackage{tikz}
\usepackage{pgfplots}
\usepackage{wasysym}
\usepackage{listings}
\usepackage{bbm}
\usepackage{hyperref}
\usepackage{enumitem}
\usepackage{transparent}
\usepackage{float}
\usepackage{multirow}
\usepackage[compat=1.1.0]{tikz-feynman}
\usepackage{slashed}
\usepackage[normalem]{ulem}
\theoremstyle{definition}
\newtheorem{Theorem}{Theorem}
\newtheorem{Proposition}{Proposition}
\newtheorem{Lemma}[Theorem]{Lemma}
\newtheorem{Corollary}[Theorem]{Corollary}
\newtheorem{Example}[Theorem]{Example}
\theoremstyle{definition}
\newtheorem{Remark}[Theorem]{Remark}
\theoremstyle{definition}
\newtheorem{Problem}{Problem}
\theoremstyle{definition}
\newtheorem{Definition}[Theorem]{Definition}
\newenvironment{parts}{\begin{enumerate}[label=(\alph*)]}{\end{enumerate}}
\newcommand{\backvec}[1]{\reflectbox{$\vec{\reflectbox{\!$#1$}}$}}

%tikz
\usetikzlibrary{patterns}
\usepackage{tikz-cd}
% definitions of number sets
\newcommand{\N}{\mathbb{N}}
\newcommand{\R}{\mathbb{R}}
\newcommand{\Z}{\mathbb{Z}}
\newcommand{\Q}{\mathbb{Q}}
\newcommand{\C}{\mathbb{C}}
\allowdisplaybreaks

\newtcolorbox{revision}[1][]{empty, breakable, sharp corners, borderline north={1pt}{0pt}{black},
	borderline west={5pt}{0pt}{black},
	#1}

\definecolor{codegreen}{rgb}{0,0.6,0}
\definecolor{codegray}{rgb}{0.5,0.5,0.5}
\definecolor{codepurple}{rgb}{0.58,0,0.82}
\definecolor{backcolour}{rgb}{0.95,0.95,0.92}

\lstdefinestyle{mystyle}{
	backgroundcolor=\color{backcolour},   
	commentstyle=\color{codegreen},
	keywordstyle=\color{magenta},
	numberstyle=\tiny\color{codegray},
	stringstyle=\color{codepurple},
	basicstyle=\ttfamily\footnotesize,
	breakatwhitespace=false,         
	breaklines=true,                 
	captionpos=b,                    
	keepspaces=true,                 
	numbers=left,                    
	numbersep=5pt,                  
	showspaces=false,                
	showstringspaces=false,
	showtabs=false,                  
	tabsize=2,
	literate={_}{{\_}}1
}

\lstset{style=mystyle}

\begin{document}
\title{Tensor Products on Quantum Spin Systems}
	\author{Jun Wei Tan}
	\email{jun-wei.tan@stud-mail.uni-wuerzburg.de}
	\affiliation{Julius-Maximilians-Universit\"{a}t W\"{u}rzburg}
	\date{\today}
	\maketitle
	
	\begin{tcolorbox}
		In the 1960s Friedrichs met Heisenberg, and used the occasion to express to him the deep gratitude of the community of mathematicians for having created quantum mechanics, which gave birth to the beautiful theory of operators in Hilbert space. Heisenberg allowed that this was so; Friedrichs then added that the mathematicians have, in some measure, returned the favor. Heisenberg looked noncommittal, so Friedrichs pointed out that it was a mathematician, von Neumann, who clarified the difference between a self-adjoint operator and one that is merely symmetric. ``What's the difference,'' said Heisenberg.
	\end{tcolorbox}
\section{Introduction - What is a Physical System?}
The question of what a physical system is goes back to the initial \textcolor{red}{DO BETTER}. Essentially, we require a $C^*$ algebra, which shall describe the possible operators that can act on a system. We call the self adjoint elements of this $C^*$ algebra the observables. To connect this to the standard Hilbert space representation taught in a first course in quantum mechanics, we can look at \emph{representations} of the $C^*$ algebra. We already covered most of these 

However, some questions remain to be answered. Most importantly, it is still not clear how to define the interactions. It is well known that the Hamiltonian is almost always unbounded; for example, the Hamiltonian of the Heisenberg chain
\[
\hat{H} = J\sum_{\langle x,y\rangle} \vec{S}_x \cdot \vec{S}_y
\] 
is unbounded on infinite systems. Thus, we often seek to define a limit in which it converges, and 

In this talk, we will work only with spin systems. A single spin is defined by the $C^*$ algebra of $2\times 2$ matrices spanned by the Pauli matrices. This is very much finite dimensional. We can hence also answer questions about two spins by finite dimensional linear algebra. In fact, we can answer questions about \emph{any finite number} of spins simply by applying finite dimensional linear algebra. Thus, by avoiding complications arising from the infinite dimensionality of individual particles, we can focus on the most interesting part of statistical physics, which is the thermodynamic limit.
\section{Dynamics on a Hilbert Space}
Previously, we defined three different notions that each purported to generate dynamics on a $C^*$ algebra $\mathcal{A}$. These were the Hamiltonian $H$, the propagator $U(t)$ and the time evolution $\alpha: \mathcal{A}\times \R \to \mathcal{A}$. We will first summarise these, before we discuss their limitations.

\begin{Definition}
	A propagator is a strongly continuous 1 parameter unitary group - that is, a function $U: \R \to \mathcal{A}$ such that $U(t)$ is unitary for each $t$ and the $U(t)$ is a group homomorphism.
\end{Definition}
\begin{Definition}
	A time evolution is a 1 parameter group of automorphisms - that is, a map $\alpha: \mathcal{A}\times \R \to \mathcal{A}$ such that $\alpha_t$ is an automorphism of $\mathcal{A}$ for each $t$.
\end{Definition}
Given a Hamiltonian $H$, we can easily define a propagator via
\[
	U(t) = e^{itH}
.\] 
Now that we have a propagator, we can define a time evolution through
\[
\alpha_t(A) = U(t)AU(t)^\dagger
.\] 
However, it is not easy to see if this can be done in reverse; in the lecture, we applied Stone's theorem to construct a Hamiltonian from a propagator	
\begin{Theorem}[Stone]
		Let $t\mapsto U(t)$ be a 1 parameter strongly continuous semigroup of unitary operators acting on some Hilbert space $\mathcal{H}$. Then there exists some a unique (possibly unbounded) operator $H$ such that
		\[U(t) = e^{itH}.\]
	\end{Theorem} 
	Thus, only one step remains to us now --- given a time evolution, how do we construct a propagator? We will see that this step has been left undone for good reason; it requires additional, highly restrictive conditions.
\begin{Theorem}[Every Automorphism of $\mathcal{B}(\mathcal{H})$ is Inner]
Consider the $C^*$ algebra $\mathcal{B}(\mathcal{H})$. Then there is a unitary operator $U\in \mathcal{B}(\mathcal{H})$ such that $\alpha(A) = UAU^*$ for all $A\in \mathcal{B}(\mathcal{H})$.
\end{Theorem}
\begin{proof}
	We do this by
\end{proof}
From the proof, we see the main issue --- here, it is absolutely imperative that we consider the $C^*$ algebra $\mathcal{B}(\mathcal{H})$, as this allowed us to make sense of 1 dimensional projectors. In the general $C^*$ algebra setting, these characteristics are all absent, defeating the purpose of an algebraic approach to quantum mechanics. Thus, we must seek a different way to construct this.

What tools do we have? In reality, the $C^*$ algebra does not contain all elements of $\mathcal{B}(\mathcal{H})$. The algebra must be \emph{local}. We will define this term rigorously later, but physically, this means that operators cannot hybridise states that are supported arbitrarily far away from each other. This leads us to our first point, the definition of the quasi-local algebra.
\section{The Quasi-Local Algebra}
Let us consider a 
\section{Dynamics of a Physical System}
\subsection{What Are Dynamics?}
In quantum mechanics, we implement the dynamics of a quantum system in the Heisenberg representation, where all observables have the time dependence
	\[A(t) = U(t)^*AU(t)\]
	where $U(t)=e^{-iHt}$ is the unitary time evolution operator generated by the Hamiltonian $H$. However, this leaves a major question. We know that $H$ is, in general, unbounded. Thus, we should not expect $U(t)$ to be well defined, or a member of the Hilbert space at all. Can we rescue it?
	
	First, let us suppose that a system already has dynamics defined on it. We shall define this as follows:
	\begin{Definition}
		A $C^*$ dynamical system is a $C^*$ algebra $\mathcal{A}$ together with a strongly continuous 1 parameter semigroup $t\mapsto U(t)$, inducing an isomorphism $\alpha_t:\mathcal{A}\to\mathcal{A}$ with $\alpha_t: A\mapsto U(t)^*A U$.  
	\end{Definition}
Now, if we were able to write this automorphism in terms of a unitary operator $U(t)$ such that $\alpha_t(A) = U(t)^* A U(t)$, we would be able to apply Stone's theorem:

We do this, in essence, by writing the automorphism in the Schrödinger picture.
\begin{Theorem}
	Let $\mathcal{A}$ be a $C^*$ subalgebra of $\mathcal{B}(\mathcal{H})$, and $\alpha:\mathcal{A}\to\mathcal{A}$ an automorphism of $\mathcal{A}$. Then there is a unitary $U$ such that $\alpha(A)=UAU^*$ for all $A\in \mathcal{A}$.
\end{Theorem}
\begin{proof}
	We define the unitary $U$ by $U:\mathcal{A}\mathcal{H}\to\mathcal{H}$ by $UA\ket{\psi}=\alpha(A)\ket{\psi'}$, where $\ket{\psi}$ is defined by $\alpha(\dyad{\psi}{\psi})=\dyad{\psi'}{\psi'}$.
	
	First, we note that $U$ is linear, which follows by using the automorphism properties of $\mathcal{B}(\mathcal{H})$ to pull out sums on the right hand side of the equation. Thus, $U$ is well defined.
	
	We then prove that $U$ is unitary. Let $\ket{\phi},\ket{\psi}\in \mathcal{H}$. We consider the subspace of operators $\dyad{\phi}{\phi}$ and $\ket{\psi}{\psi}$. Since these are rank 1 operators, the proof transfers over from finite dimensional linear algebra: We have
	\[\tr(\ket{\psi}\braket{\psi}{\phi}\bra{\phi})=|\braket{\psi}{\phi}|^2\]
	Since automorphisms preserve the trace, we have
	\[|\braket{\psi}{\phi}|^2 = |\braket{\psi'}{\phi'}|^2 = |\mel{\psi}{U^*U}{\phi}|^2\]
	and thus $U$ is unitary. Finally, since \textcolor{red}{TODO}
\end{proof}

\end{document}
